
\documentclass[11pt]{article}

\usepackage[margin=1in]{geometry}
\usepackage{amsmath,amssymb,amsthm,mathtools}
\usepackage{bm}
\usepackage{enumitem}
\usepackage{xcolor}
\usepackage{hyperref}

\newcommand{\RR}{\mathbb{R}}
\newcommand{\rank}{\operatorname{rank}}
\newcommand{\range}{\operatorname{range}}
\newcommand{\col}{\operatorname{col}}
\newcommand{\tr}{\operatorname{tr}}
\newcommand{\diag}{\operatorname{diag}}
\newcommand{\Span}{\operatorname{span}}
\newcommand{\cD}{\mathcal{D}}
\newcommand{\cE}{\mathcal{E}}
\newcommand{\cS}{\mathcal{S}}
\newcommand{\cK}{\mathcal{K}}
\newcommand{\cU}{\mathcal{U}}
\newcommand{\sol}[1]{\textcolor{red}{\textbf{[Sol: #1]}}}

\newtheorem{theorem}{Theorem}
\newtheorem{lemma}{Lemma}
\newtheorem{proposition}{Proposition}
\newtheorem{corollary}{Corollary}
\newtheorem{remark}{Remark}
\newtheorem{definition}{Definition}

\title{Almost-Global Convergence of Preconditioned Gradient Descent for Low-Rank PSD Approximation}
\author{}
\date{}

\begin{document}

\maketitle

\section{Problem setup}

Let
\[
M\in\RR^{n\times n}
\]
be symmetric, and consider
\[
g(X)
=
\frac14\left\|M-XX^\top\right\|_F^2,
\qquad
X\in\RR^{n\times k},
\]
where \(k\le n\). We study the preconditioned gradient descent iteration
\[
X_{t+1}
=
X_t
-
\eta \nabla g(X_t)(X_t^\top X_t)^{-1}.
\]

The desired conclusion cannot hold for an arbitrary indefinite symmetric matrix.
We therefore assume throughout that
\[
M\succeq 0,
\qquad
k\le \rank(M).
\]
Under these assumptions, minimizing \(g\) recovers the rank-\(k\) truncated
eigendecomposition of \(M\), equivalently its rank-\(k\) singular value
decomposition.

Let
\[
\cS:=\range(M),
\qquad
r:=\rank(M),
\]
and define the full-rank support domain
\[
\cD_M
:=
\left\{
X\in \cS^{\times k}:\rank(X)=k
\right\}.
\]
For \(X\in\cD_M\), define
\[
\Phi(X)
:=
\frac12\tr(X^\top M^\dagger X)
-
\frac12\log\det(X^\top X).
\]

\section{Main results}

\begin{theorem}[Uniform almost-global convergence on a potential sublevel]
\label{thm:main}
Suppose that \(M\succeq0\) has rank \(r\ge k\), with positive spectral
decomposition
\[
M
=
U_r\Lambda U_r^\top,
\qquad
\Lambda
=
\diag(\lambda_1,\ldots,\lambda_r),
\qquad
\lambda_1\ge\cdots\ge\lambda_r>0.
\]
For \(C\in\RR\), let
\[
\cK_C^X
:=
\left\{
X\in\cD_M:\Phi(X)\le C
\right\},
\]
and assume \(\cK_C^X\neq\varnothing\).

There exists
\[
\eta_C=\eta_C(M,k,C)>0
\]
such that, for every fixed \(0<\eta<\eta_C\), the following statements hold.

\begin{enumerate}[label=(\roman*)]
    \item For every \(X_0\in\cK_C^X\), all iterates are well-defined,
    remain full column rank, stay in \(\cK_C^X\), and converge to a critical
    point of \(g\).

    \item For Lebesgue-almost every \(X_0\in\cK_C^X\), the limit is a global
    minimizer. More precisely,
    \[
    X_t\longrightarrow
    X_\infty
    =
    U_\star\Lambda_\star^{1/2}R,
    \qquad
    R\in O(k),
    \]
    where \(U_\star\) spans a top-\(k\) invariant subspace of \(M\). Hence
    \[
    X_\infty X_\infty^\top
    =
    U_\star\Lambda_\star U_\star^\top
    \]
    is a best rank-\(k\) approximation of \(M\), and
    \[
    g(X_\infty)
    =
    \frac14\sum_{i=k+1}^r\lambda_i^2.
    \]
\end{enumerate}

Here, ``almost every'' is with respect to the \(rk\)-dimensional Lebesgue
measure on
\[
\cS^{\times k}\cong\RR^{r\times k}.
\]
If \(M\succ0\), this is the ordinary ambient Lebesgue measure on
\(\RR^{n\times k}\). If, additionally,
\[
\lambda_k>\lambda_{k+1},
\qquad
\lambda_{r+1}:=0,
\]
then
\[
\col(X_t)
\longrightarrow
\Span(u_1,\ldots,u_k).
\]
\end{theorem}

\sol{The theorem is uniform over a prescribed potential sublevel. The proof does not establish one initialization-independent fixed step size on the entire unbounded full-rank domain.}

\begin{corollary}[Almost every initialization with an energy-selected fixed step]
\label{cor:shell}
For every integer \(j\) for which
\[
\cE_j
:=
\left\{
X\in\cD_M:j-1<\Phi(X)\le j
\right\}
\]
is nonempty, let \(\overline\eta_j>0\) be any valid threshold supplied by
Theorem~\ref{thm:main} with \(C=j\), and choose once and for all
\[
0<\eta_j<\overline\eta_j.
\]

Given \(X_0\in\cE_j\), run the iteration with the single fixed step size
\(\eta=\eta_j\). Then the algorithm converges to a global minimizer for
Lebesgue-almost every \(X_0\in\cD_M\).
\end{corollary}

\begin{proof}
On each shell \(\cE_j\), the exceptional set has measure zero by
Theorem~\ref{thm:main}. The domain \(\cD_M\) is the countable union of the
shells \(\cE_j\), and a countable union of measure-zero sets has measure zero.
\end{proof}

\sol{When \(M\) is singular, this is a relative-measure statement on \(\range(M)^{\times k}\), not an ambient almost-everywhere statement in \(\RR^{n\times k}\). A practical random initialization is given in Section~\ref{sec:singular-init}.}

\section{Rewriting the iteration}

We first compute the gradient of \(g\).

\begin{lemma}
For every \(X\in\RR^{n\times k}\),
\[
\nabla g(X)
=
(XX^\top-M)X
=
X(X^\top X)-MX.
\]
\end{lemma}

\begin{proof}
Expanding the objective gives
\[
g(X)
=
\frac14\tr(M^2)
-
\frac12\tr(X^\top MX)
+
\frac14\tr\bigl((X^\top X)^2\bigr).
\]
Because \(M\) is symmetric,
\[
\nabla_X\left(-\frac12\tr(X^\top MX)\right)
=
-MX,
\]
while
\[
\nabla_X\left(\frac14\tr\bigl((X^\top X)^2\bigr)\right)
=
X(X^\top X).
\]
Combining the two identities proves the claim.
\end{proof}

It follows that
\[
\nabla g(X)(X^\top X)^{-1}
=
X-MX(X^\top X)^{-1},
\]
so the iteration may be written as
\[
X_{t+1}
=
(1-\eta)X_t
+
\eta MX_t(X_t^\top X_t)^{-1}.
\tag{1}
\label{eq:x-update}
\]
Equation~\eqref{eq:x-update} shows that if the columns of \(X_0\) lie in
\(\range(M)\), then the same is true for every \(X_t\).

Recall the positive spectral decomposition
\[
M=U_r\Lambda U_r^\top,
\qquad
\Lambda=\diag(\lambda_1,\ldots,\lambda_r)\succ0.
\]
Since \(X_t\in\range(M)^{\times k}\), there is a unique
\(\overline X_t\in\RR^{r\times k}\) such that
\[
X_t=U_r\overline X_t,
\qquad
\overline X_t=U_r^\top X_t.
\]
Because \(U_r^\top U_r=I_r\), the objective, potential, and iteration reduce to
\[
g(U_r\overline X)
=
\frac14
\left\|
\Lambda-\overline X\overline X^\top
\right\|_F^2,
\]
\[
\Phi(U_r\overline X)
=
\frac12\tr(\overline X^\top\Lambda^{-1}\overline X)
-
\frac12\log\det(\overline X^\top\overline X),
\]
and
\[
\overline X_{t+1}
=
(1-\eta)\overline X_t
+
\eta\Lambda\overline X_t
(\overline X_t^\top\overline X_t)^{-1}.
\]
Here \(\Lambda\in\RR^{r\times r}\) is positive definite even when the original
ambient matrix \(M\in\RR^{n\times n}\) is singular. From this point until
Section~\ref{sec:translation-back}, the proof is carried out entirely in the
reduced coordinates \(\RR^{r\times k}\), and all inverses and inverse square
roots are applied to \(\Lambda\), not to the possibly singular ambient matrix
\(M\).

\section{A hidden log-determinant potential}

On
\[
\cD
:=
\left\{
\overline X\in\RR^{r\times k}:\rank(\overline X)=k
\right\},
\]
define the reduced potential
\[
\overline\Phi(\overline X)
:=
\frac12\tr(\overline X^\top\Lambda^{-1}\overline X)
-
\frac12\log\det(\overline X^\top\overline X).
\]
By the preceding reduction,
\[
\overline\Phi(\overline X)=\Phi(U_r\overline X).
\]
Its Euclidean gradient is
\[
\nabla\overline\Phi(\overline X)
=
\Lambda^{-1}\overline X
-
\overline X(\overline X^\top\overline X)^{-1},
\]
and therefore
\[
\Lambda\nabla\overline\Phi(\overline X)
=
\overline X
-
\Lambda\overline X(\overline X^\top\overline X)^{-1}.
\]
Consequently, the reduced iteration is gradient descent on
\(\overline\Phi\) with respect to the fixed inner product
\[
\langle A,B\rangle_{\Lambda^{-1}}
:=
\tr(A^\top\Lambda^{-1}B).
\]

For an ordinary Euclidean formulation, set
\[
\overline X=\Lambda^{1/2}Y,
\qquad
Y=\Lambda^{-1/2}\overline X,
\]
and define
\[
F(Y)
:=
\overline\Phi(\Lambda^{1/2}Y)
=
\frac12\|Y\|_F^2
-
\frac12\log\det(Y^\top\Lambda Y).
\tag{2}
\label{eq:F-def}
\]
Let
\[
A(Y):=Y^\top\Lambda Y.
\]
Differentiating gives
\[
\nabla F(Y)
=
Y-\Lambda YA(Y)^{-1}.
\tag{3}
\label{eq:F-grad}
\]
Substituting \(\overline X=\Lambda^{1/2}Y\) into the reduced iteration yields
\[
Y_{t+1}
=
Y_t-\eta\nabla F(Y_t).
\tag{4}
\label{eq:Y-update}
\]
Thus the preconditioned method in the reduced variable \(\overline X\) is
exactly ordinary Euclidean gradient descent on the real-analytic function
\(F\).

\begin{lemma}[Full-rank preservation]
\label{lem:rank-preservation}
If \(Y\) has full column rank and \(0<\eta\le1\), then
\[
Y_+
:=(1-\eta)Y+\eta\Lambda Y(Y^\top\Lambda Y)^{-1}
\]
also has full column rank.
\end{lemma}

\begin{proof}
We have
\[
Y^\top Y_+
=
(1-\eta)Y^\top Y+\eta I_k\succ0.
\]
If \(Y_+z=0\), then
\[
0
=
z^\top Y^\top Y_+z
=
z^\top\bigl((1-\eta)Y^\top Y+\eta I_k\bigr)z,
\]
which forces \(z=0\). Therefore \(Y_+\) has trivial kernel.
\end{proof}

\section{Compactness of sublevel sets}

For \(C\in\RR\), define
\[
\cK_C
:=
\left\{
Y\in\cD:F(Y)\le C
\right\}.
\]

\begin{lemma}[Coercivity and a barrier at rank deficiency]
\label{lem:coercivity}
Every nonempty sublevel set \(\cK_C\) is compact in
\(\RR^{r\times k}\) and is contained a positive distance from the rank-deficient
boundary of \(\cD\).
\end{lemma}

\begin{proof}
If a bounded sequence \(Y_m\in\cD\) approaches rank deficiency, then
\[
\det(Y_m^\top\Lambda Y_m)\longrightarrow0,
\]
and hence \(F(Y_m)\to+\infty\).

It remains to rule out escape to infinity. Let
\[
L:=\lambda_1(\Lambda)=\lambda_1(M),
\qquad
s:=\|Y\|_F^2.
\]
By the arithmetic--geometric mean inequality,
\[
\det(Y^\top\Lambda Y)
\le
\left(
\frac{\tr(Y^\top\Lambda Y)}{k}
\right)^k
\le
\left(\frac{Ls}{k}\right)^k.
\]
Therefore
\[
F(Y)
\ge
\frac{s}{2}
-
\frac{k}{2}\log\left(\frac{Ls}{k}\right),
\tag{5}
\label{eq:coercive-lower}
\]
and the right-hand side tends to \(+\infty\) as \(s\to\infty\).

Thus \(\cK_C\) is bounded. If \(Y_m\in\cK_C\) converges in the ambient
Euclidean space, the barrier argument prevents its limit from being rank
deficient; continuity of \(F\) then shows that the limit remains in \(\cK_C\).
Hence \(\cK_C\) is closed and bounded, and therefore compact.
\end{proof}

In particular, there exist constants
\[
0<\delta_C<R_C<\infty
\]
such that
\[
\sigma_{\min}(Y)\ge\delta_C,
\qquad
\|Y\|_F\le R_C
\]
for every \(Y\in\cK_C\). This replaces any need for a translated cone, slab,
or local-neighborhood assumption.

\section{A sufficient step-size bound}

Assume \(\cK_C\neq\varnothing\). Define
\[
S_C
:=
\max\left\{
s>0:
\frac{s}{2}
-
\frac{k}{2}\log\left(\frac{Ls}{k}\right)
\le C
\right\}.
\tag{6}
\label{eq:S-C}
\]
By \eqref{eq:coercive-lower},
\[
\|Y\|_F^2\le S_C
\]
for every \(Y\in\cK_C\).

From
\[
F(Y)
=
\frac12\|Y\|_F^2
-
\frac12\log\det A(Y),
\]
we obtain
\[
\det A(Y)
=
\exp\bigl(\|Y\|_F^2-2F(Y)\bigr)
\ge
e^{-2C}.
\]
Also,
\[
\lambda_{\max}(A(Y))
\le
L\|Y\|_F^2
\le
LS_C.
\]
Therefore
\[
\lambda_{\min}(A(Y))
\ge
a_C
:=
\frac{e^{-2C}}{(LS_C)^{k-1}}.
\tag{7}
\label{eq:a-C}
\]

Using \eqref{eq:F-grad}, on \(\cK_C\) we have
\[
\|\nabla F(Y)\|_F
\le
\sqrt{S_C}
\left(
1+\frac{L}{a_C}
\right)
=:
\Gamma_C.
\tag{8}
\label{eq:Gamma-C}
\]
The Hessian action is
\[
D\nabla F(Y)[H]
=
H-\Lambda H A^{-1}
+
\Lambda Y A^{-1}
\left(
H^\top\Lambda Y+Y^\top\Lambda H
\right)
A^{-1},
\tag{9}
\label{eq:F-Hessian}
\]
where \(A=A(Y)\). Whenever
\[
\|Y\|_F<\sqrt{S_C}+1,
\qquad
\lambda_{\min}(A(Y))>\frac{a_C}{2},
\]
we may take
\[
\|D\nabla F(Y)\|_{\mathrm{op}}
\le
\mathcal L_C,
\]
where
\[
\mathcal L_C
:=
1
+
\frac{2L}{a_C}
+
\frac{8L^2(\sqrt{S_C}+1)^2}{a_C^2}.
\tag{10}
\label{eq:L-C}
\]

Define
\[
d_C
:=
\min\left\{
1,
\frac{a_C}{4L(\sqrt{S_C}+1)}
\right\},
\tag{11}
\label{eq:d-C}
\]
and choose
\[
0<\eta<\eta_C,
\qquad
\eta_C
:=
\min\left\{
1,
\frac{d_C}{\Gamma_C},
\frac{1}{\mathcal L_C}
\right\}.
\tag{12}
\label{eq:eta-C}
\]

Let
\[
\Delta=-\eta\nabla F(Y).
\]
For \(Y\in\cK_C\), \eqref{eq:Gamma-C} gives
\[
\|\Delta\|_F<d_C.
\]
For every \(s\in[0,1]\),
\[
A(Y+s\Delta)-A(Y)
=
s\Delta^\top\Lambda Y
+
sY^\top\Lambda\Delta
+
s^2\Delta^\top\Lambda\Delta.
\]
Hence
\[
\|A(Y+s\Delta)-A(Y)\|_2
\le
2L\sqrt{S_C}\,d_C+Ld_C^2
\le
\frac{a_C}{2}.
\]
The entire update segment remains full rank and lies in the region where the
Hessian is bounded by \(\mathcal L_C\). Taylor's theorem therefore gives
\[
\begin{aligned}
F(Y-\eta\nabla F(Y))
&\le
F(Y)
-
\eta\|\nabla F(Y)\|_F^2
+
\frac{\mathcal L_C\eta^2}{2}\|\nabla F(Y)\|_F^2\\
&\le
F(Y)
-
\frac{\eta}{2}\|\nabla F(Y)\|_F^2.
\end{aligned}
\tag{13}
\label{eq:descent}
\]
Consequently,
\[
F(Y_{t+1})\le F(Y_t)\le F(Y_0)\le C,
\]
so every iterate remains in \(\cK_C\).

\section{Convergence to a critical point}

Summing \eqref{eq:descent} gives
\[
\frac{\eta}{2}
\sum_{t=0}^\infty
\|\nabla F(Y_t)\|_F^2
\le
F(Y_0)-\inf F
<\infty.
\]
Therefore
\[
\|\nabla F(Y_t)\|_F\longrightarrow0.
\tag{14}
\label{eq:grad-to-zero}
\]
Moreover,
\[
Y_{t+1}-Y_t
=
-\eta\nabla F(Y_t),
\tag{15}
\label{eq:relative-error}
\]
and \eqref{eq:descent} can be rewritten as
\[
F(Y_t)-F(Y_{t+1})
\ge
\frac{1}{2\eta}\|Y_{t+1}-Y_t\|_F^2.
\tag{16}
\label{eq:strong-descent}
\]

The sequence lies in the compact set \(\cK_C\), and every cluster point is
critical by \eqref{eq:grad-to-zero}. Since \(F\) is real analytic on the
full-rank domain, the discrete \L{}ojasiewicz convergence theorem, applied
using \eqref{eq:relative-error} and \eqref{eq:strong-descent}, yields
\[
\sum_{t=0}^\infty
\|Y_{t+1}-Y_t\|_F
<
\infty.
\]
Thus \(Y_t\) converges to a single critical point.

\sol{The previous version invoked \L{}ojasiewicz convergence without displaying the required descent and relative-error hypotheses. They are now explicit; Appendix~\ref{app:loj} records the precise external theorem being used.}

\section{Characterization of all critical points}

\begin{proposition}
\label{prop:critical-points}
A full-column-rank matrix \(Y\) is a critical point of \(F\) if and only if
\begin{enumerate}[label=(\roman*)]
    \item \(Y^\top Y=I_k\), and
    \item \(\col(Y)\) is an invariant subspace of \(\Lambda\).
\end{enumerate}
\end{proposition}

\begin{proof}
At a critical point,
\[
Y-\Lambda Y(Y^\top\Lambda Y)^{-1}=0.
\]
Let
\[
A:=Y^\top\Lambda Y.
\]
Then
\[
\Lambda Y=YA.
\tag{17}
\label{eq:invariance}
\]
Multiplying the critical-point equation on the left by \(Y^\top\) gives
\[
Y^\top Y
=
Y^\top\Lambda Y\,A^{-1}
=
I_k.
\tag{18}
\label{eq:orthonormal}
\]
Thus \(Y\) has orthonormal columns and \(\col(Y)\) is invariant under \(\Lambda\).

Conversely, suppose that \(Y^\top Y=I_k\) and that \(\Lambda Y=YA\) for some
symmetric \(A\in\RR^{k\times k}\). Then \(Y^\top\Lambda Y=A\), and hence
\[
\Lambda Y(Y^\top\Lambda Y)^{-1}
=
YAA^{-1}
=
Y.
\]
Therefore \(\nabla F(Y)=0\).
\end{proof}

\begin{lemma}[Right-orthogonal eigenbasis reduction]
\label{lem:eigenbasis-rotation}
Let \(Y\) be a critical point. There exists \(R\in O(k)\) such that the
columns of \(YR\) are eigenvectors of \(\Lambda\). Moreover,
\[
F(YR)=F(Y),
\]
and the Hessians at \(Y\) and \(YR\) have the same inertia.
\end{lemma}

\begin{proof}
At a critical point, \(\Lambda Y=YA\), where \(A=Y^\top\Lambda Y\) is symmetric positive
definite. Choose \(R\in O(k)\) such that
\[
R^\top AR=\diag(\mu_1,\ldots,\mu_k).
\]
Then
\[
\Lambda(YR)
=
YR\diag(\mu_1,\ldots,\mu_k),
\]
so the columns of \(YR\) are eigenvectors. Right-orthogonal invariance follows
from
\[
\|YR\|_F=\|Y\|_F,
\qquad
\det(R^\top Y^\top\Lambda YR)=\det(Y^\top\Lambda Y).
\]
Because the map \(H\mapsto HR\) is an isometry and \(F(\cdot R)=F(\cdot)\),
the Hessian quadratic forms are related by this isometry and therefore have
the same inertia.
\end{proof}

Suppose that a critical invariant subspace corresponds to eigenvalues
\[
\mu_1,\ldots,\mu_k.
\]
Then
\[
F(Y)
=
\frac{k}{2}
-
\frac12\sum_{i=1}^k\log\mu_i.
\tag{19}
\label{eq:critical-value}
\]
Therefore the global minima are precisely the invariant subspaces selecting
the \(k\) largest eigenvalues of \(\Lambda\).

\section{Nonglobal critical points are strict saddles}

\begin{proposition}
\label{prop:strict-saddle}
Every nonglobal critical point of \(F\) has a direction of strictly negative
curvature.
\end{proposition}

\begin{proof}
Let \(Y\) be a nonglobal critical point. By
Lemma~\ref{lem:eigenbasis-rotation}, after a right-orthogonal rotation we may
assume that the columns of \(Y\) are orthonormal eigenvectors of \(\Lambda\).

Because the selected invariant subspace is not globally optimal, there exists
a selected eigenvector \(u\in\col(Y)\) and an unselected eigenvector
\(v\perp\col(Y)\) such that
\[
\Lambda u=\alpha u,
\qquad
\Lambda v=\beta v,
\qquad
\beta>\alpha.
\]
Assume that \(u\) is the \(i\)-th column of \(Y\), and define
\[
H:=ve_i^\top.
\]
The second variation of \(F\) is
\[
\begin{aligned}
\nabla^2F(Y)[H,H]
={}&
\|H\|_F^2
-
\tr\left(A^{-1}H^\top\Lambda H\right)\\
&+
\frac12
\tr\left[
A^{-1}
\left(H^\top\Lambda Y+Y^\top\Lambda H\right)
A^{-1}
\left(H^\top\Lambda Y+Y^\top\Lambda H\right)
\right],
\end{aligned}
\tag{20}
\label{eq:second-variation}
\]
where \(A=Y^\top\Lambda Y\). Since \(v\perp\col(Y)\),
\[
H^\top\Lambda Y=0,
\qquad
Y^\top\Lambda H=0.
\]
Also,
\[
H^\top\Lambda H=\beta e_ie_i^\top,
\]
and the \(i\)-th diagonal entry of \(A\) is \(\alpha\). Therefore
\[
\nabla^2F(Y)[H,H]
=
1-\frac{\beta}{\alpha}
<0.
\]
Hence \(Y\) is a strict saddle.
\end{proof}

No eigengap assumption is required for Proposition~\ref{prop:strict-saddle}.
If the \(k\)-th eigenvalue is repeated, every admissible \(k\)-dimensional
choice within the tied top eigenspace is itself a global minimizer.

\section{Almost-sure avoidance of strict saddles}

Let
\[
T_\eta(Y)
:=
Y-\eta\nabla F(Y).
\]
The compact sublevel \(\cK_C\) lies in the open set
\[
\cU_C
:=
\left\{
Y:\|Y\|_F<\sqrt{S_C}+1,
\ \lambda_{\min}(Y^\top\Lambda Y)>\frac{a_C}{2}
\right\},
\]
on which
\[
\|\nabla^2F(Y)\|_{\mathrm{op}}\le\mathcal L_C.
\]
Since \(\eta\mathcal L_C<1\),
\[
DT_\eta(Y)
=
I-\eta\nabla^2F(Y)
\]
is nonsingular throughout \(\cU_C\). Hence \(T_\eta\) is a local
\(C^1\)-diffeomorphism there.

At a strict saddle \(Y_s\), the Hessian has a negative eigenvalue
\(-\gamma\), so \(DT_\eta(Y_s)\) has an eigenvalue
\[
1+\eta\gamma>1.
\]
The local center-stable manifold theorem therefore places all points whose
forward trajectories remain sufficiently near \(Y_s\) in a manifold of
dimension at most \(rk-1\), hence in a Lebesgue-null set. Using a countable
cover of the strict-saddle set and the fact that local diffeomorphisms preserve
nullity under inverse images, the basin of all strict saddles inside
\(\cK_C\) is a countable union of null sets and is therefore null.

Since every trajectory in \(\cK_C\) converges to a critical point and every
nonglobal critical point is a strict saddle, almost every initialization in
\(\cK_C\) converges to a global minimum of \(F\).

\sol{The critical sets are generally non-isolated because \(Y\mapsto YR\) leaves \(F\) unchanged. Thus one must use the non-isolated center-stable argument, not only the simpler isolated-saddle theorem. Full details are in Appendix~\ref{app:saddles}.}

\section{Translation back to the original objective}
\label{sec:translation-back}

The preceding sections analyzed the reduced variables
\[
X_t=U_r\overline X_t,
\qquad
\overline X_t=\Lambda^{1/2}Y_t.
\]
At a global minimum of the reduced potential,
\[
Y_\infty
=
E_\star R,
\qquad
R\in O(k),
\]
where \(E_\star\in\RR^{r\times k}\) has orthonormal columns spanning a
top-\(k\) invariant subspace of \(\Lambda\). Writing
\[
\Lambda E_\star
=
E_\star\Lambda_\star,
\]
we obtain
\[
\overline X_\infty
=
\Lambda^{1/2}Y_\infty
=
E_\star\Lambda_\star^{1/2}R.
\]
Returning to the original ambient coordinates \(\RR^n\) gives
\[
\begin{aligned}
X_\infty
&=
U_r\overline X_\infty\\
&=
U_rE_\star\Lambda_\star^{1/2}R\\
&=
U_\star\Lambda_\star^{1/2}R,
\end{aligned}
\]
where
\[
U_\star:=U_rE_\star\in\RR^{n\times k}.
\]
Since
\[
M U_\star
=
U_r\Lambda E_\star
=
U_rE_\star\Lambda_\star
=
U_\star\Lambda_\star,
\]
the columns of \(U_\star\) span a top-\(k\) invariant subspace of the original
ambient matrix \(M\). Consequently,
\[
X_\infty X_\infty^\top
=
U_\star\Lambda_\star U_\star^\top.
\tag{21}
\label{eq:truncated-M}
\]
By the Eckart--Young--Mirsky theorem,
\[
\min_{\rank(Z)\le k}\|M-Z\|_F^2
=
\sum_{i=k+1}^r\lambda_i^2.
\]
The optimal rank-\(k\) approximation is positive semidefinite and therefore
can be written as \(XX^\top\). Thus
\[
\min_X g(X)
=
\frac14\sum_{i=k+1}^r\lambda_i^2,
\]
and \eqref{eq:truncated-M} is a global minimizer of \(g\).

\section{Initialization for singular \texorpdfstring{\(M\)}{M}}
\label{sec:singular-init}

If \(M\) is singular, initialization inside \(\range(M)\) does not require
knowing the leading eigenvectors. Draw
\[
\Omega\in\RR^{n\times k}
\]
from any distribution absolutely continuous with respect to ambient Lebesgue
measure, and set
\[
X_0=M\Omega.
\tag{22}
\label{eq:random-init}
\]
Since \(k\le\rank(M)\),
\[
\rank(M\Omega)=k
\]
with probability one. The induced law of \(X_0\) is absolutely continuous with
respect to the \(rk\)-dimensional Lebesgue measure on
\(\range(M)^{\times k}\). Moreover, \(X_0\in\range(M)^{\times k}\), and this
property is preserved by the iteration.

For this initialization, the initial potential can be evaluated without
forming \(M^\dagger\):
\[
\Phi(M\Omega)
=
\frac12\tr(\Omega^\top M\Omega)
-
\frac12\log\det(\Omega^\top M^2\Omega).
\]
Thus one may select a fixed step from the initial potential shell as in
Corollary~\ref{cor:shell} without computing a leading eigenspace.

\section{Conclusion}

Under the natural assumptions
\[
M\succeq0,
\qquad
k\le\rank(M),
\qquad
X_0\in\range(M)^{\times k},
\qquad
\rank(X_0)=k,
\]
preconditioned gradient descent converges to a global minimizer for almost
every initialization, provided the fixed step size is chosen uniformly on a
bounded potential sublevel, or equivalently selected once from the initial
potential shell. No translated-cone condition, slab condition, local
initialization condition, or prior approximation to the leading eigenspace is
required.

\appendix

\section{Audit notes and remaining scope}
\label{app:audit}

This appendix records the longer rigor checks underlying the short red comments
in the main text.

\subsection{What the fixed-step theorem does and does not prove}

The proved quantifiers are
\[
\forall C\quad
\exists\eta_C>0\quad
\forall\eta\in(0,\eta_C):
\quad
\text{almost every }X_0\in\cK_C^X
\text{ converges globally}.
\]
This is a uniform theorem on every prescribed potential sublevel. It does not,
by itself, imply
\[
\exists\eta>0:
\quad
\text{almost every }X_0\in\cD_M
\text{ converges globally with that same }\eta.
\]
The latter would require a global smoothness or invariant-region argument that
is uniform as \(X_0\) approaches rank deficiency or infinity. The current
proof does not provide such an argument.

Corollary~\ref{cor:shell} nevertheless covers almost every point of the entire
support domain: the step size is selected once from \(X_0\), then held fixed
throughout the run. This is not a local spectral initialization assumption,
because \(\Phi(X_0)\) controls scale and conditioning rather than angular
closeness to the leading eigenspace.

\subsection{Relative measure for singular matrices}

If \(r<n\), then \(\range(M)^{\times k}\) has dimension \(rk<nk\), so it has
ambient Lebesgue measure zero in \(\RR^{n\times k}\). Consequently, the correct
claim is almost-everywhere convergence relative to the support. The random map
\(\Omega\mapsto M\Omega\) provides a natural way to sample an absolutely
continuous distribution on that support.

\subsection{Why \texorpdfstring{\(M\succeq0\)}{positive semidefiniteness} is essential}

For \(n=k=1\) and \(M=[-1]\),
\[
g(x)=\frac14(1+x^2)^2
\]
has unique minimizer \(x=0\), but the preconditioned update is
\[
x_{t+1}
=
(1-\eta)x_t-\frac{\eta}{x_t}.
\]
No nonzero trajectory can converge to zero because the last term diverges.
Moreover, \(XX^\top\) is always positive semidefinite, so the objective does
not represent the usual truncated SVD of an indefinite symmetric matrix.

\subsection{Why the undamped step can fail}

For \(n=k=1\), \(M=[\lambda]\), and \(\eta=1\),
\[
x_{t+1}=\frac{\lambda}{x_t}.
\]
Unless \(x_0^2=\lambda\), this is an exact two-cycle. Thus damping is not merely
an artifact of the proof.

\section{The discrete \L{}ojasiewicz step}
\label{app:loj}

A standard analytic-descent theorem states that if a bounded sequence
\(z_t\) is generated for a real-analytic function \(f\) and satisfies, for
fixed constants \(a,b>0\),
\[
f(z_t)-f(z_{t+1})
\ge
a\|z_{t+1}-z_t\|^2
\]
and
\[
\|\nabla f(z_t)\|
\le
b\|z_{t+1}-z_t\|,
\]
then the sequence has finite length and converges to a single critical point.
In the present setting, \eqref{eq:strong-descent} gives
\[
a=\frac{1}{2\eta},
\]
while \eqref{eq:relative-error} gives the equality
\[
\|\nabla F(Y_t)\|_F
=
\frac1\eta\|Y_{t+1}-Y_t\|_F.
\]
The trajectory is bounded because it remains in the compact set \(\cK_C\),
and \(F\) is real analytic there. Thus all hypotheses of the theorem are
satisfied. See, for example, Absil, Mahony, and Andrews
\cite{absil2005convergence}.

\section{The non-isolated strict-saddle argument}
\label{app:saddles}

For completeness, we spell out the measure argument suppressed in the main
text.

Let \(d=rk\). For every strict saddle \(Y_s\), the derivative
\(DT_\eta(Y_s)\) has an eigenvalue of modulus strictly larger than one. The
local center-stable manifold theorem therefore provides neighborhoods
\[
V_s\Subset B_s
\]
and a local center-stable manifold \(W_s^{cs}\subset B_s\) of dimension at
most \(d-1\) such that every point whose entire forward orbit remains in
\(B_s\) lies in \(W_s^{cs}\). Hence \(W_s^{cs}\) has \(d\)-dimensional
Lebesgue measure zero.

Because \(\RR^{r\times k}\) is second countable, the strict-saddle set admits
a countable subcover \(\{V_{s_j}\}_{j\ge1}\). If an orbit converges to a strict
saddle \(Y_\infty\), then \(Y_\infty\in V_{s_j}\) for some \(j\), and for some
integer \(N\),
\[
T_\eta^N(Y_0)\in W_{s_j}^{cs}.
\]
Thus the full strict-saddle basin is contained in
\[
\bigcup_{j=1}^{\infty}
\bigcup_{N=0}^{\infty}
T_\eta^{-N}(W_{s_j}^{cs}).
\tag{23}
\label{eq:saddle-basin-cover}
\]

It remains to justify that inverse images of null sets are null. On \(\cU_C\),
\(T_\eta\) is a local \(C^1\)-diffeomorphism. Cover \(\cU_C\) by countably many
inverse-function neighborhoods. After restricting to relatively compact
subsets, each local inverse is Lipschitz, and Lipschitz maps send null sets to
null sets. Therefore
\[
T_\eta^{-1}(N)\cap\cU_C
\]
is null whenever \(N\) is null. Iterating this observation shows that every
set in \eqref{eq:saddle-basin-cover} is null. The basin is a countable union of
null sets and is therefore null.

The use of a non-isolated version is important: the right-orthogonal symmetry
\(Y\mapsto YR\) creates continuous manifolds of critical points. See
Panageas and Piliouras \cite{panageas2017nonisolated}; for the isolated case,
see Lee et al.\ \cite{lee2016gradient}.

\section{Conservativeness of the explicit step bound}

The constants in \eqref{eq:eta-C} are designed to make the update segment stay
inside an explicitly controlled full-rank neighborhood. They are sufficient,
not sharp, and may be extremely small when the sublevel permits poorly
conditioned matrices. Improving the practical step-size range is a separate
quantitative problem and is not needed for the qualitative almost-global
convergence theorem.

\begin{thebibliography}{9}

\bibitem{absil2005convergence}
P.-A. Absil, R. Mahony, and B. Andrews.
\newblock Convergence of the iterates of descent methods for analytic cost
functions.
\newblock \emph{SIAM Journal on Optimization}, 16(2):531--547, 2005.

\bibitem{lee2016gradient}
J. D. Lee, M. Simchowitz, M. I. Jordan, and B. Recht.
\newblock Gradient descent only converges to minimizers.
\newblock In \emph{Proceedings of the 29th Conference on Learning Theory},
pages 1246--1257, 2016.

\bibitem{panageas2017nonisolated}
I. Panageas and G. Piliouras.
\newblock Gradient descent only converges to minimizers: Non-isolated critical
points and invariant regions.
\newblock In \emph{8th Innovations in Theoretical Computer Science Conference},
2017.

\end{thebibliography}

\end{document}
